%% ma: L10-B01-0012
%% lop: 10
%% chuong: 1
%% bai: 1
%% dang: 10.1.5
%% loai: DS
%% muc: [NB, TH, NB, NB]
%% dap_an: "ĐSĐĐ"
%% nguon: "Ảnh giáo viên gửi 10/10/2026 (ThS. Nguyễn Hoàng Việt, Chương 1 Mệnh đề), A-Đề 1, Phần 2 Câu 13"
%% kien_thuc: "Bài 1 (mệnh đề phủ định)"
%% trang_thai: cho-duyet
%% ly_do: "Dạng toán mới đề xuất 10.1.5: mệnh đề phủ định"
\begin{ex}
Xét tính đúng, sai của các câu sau
\choiceTF
{\True $P$: ``$3^3$ là số chính phương'' có mệnh đề phủ định là $\overline P$: ``$3^3$ không là số chính phương''.}
{$Q$: ``Tam giác $ABC$ là tam giác cân'' có mệnh đề phủ định là $\overline Q$: ``Tam giác $ABC$ không là tam giác vuông''.}
{\True $R$: ``$2^{2003}-1$ là số nguyên tố'' có mệnh đề phủ định là $\overline R$: ``$2^{2003}-1$ không là số nguyên tố''.}
{\True $H$: ``$\sqrt2$ là số vô tỉ'' có mệnh đề phủ định là $\overline H$: ``$\sqrt2$ là số hữu tỉ''.}
\loigiai{a) Đúng, thêm ``không'' vào vị từ. b) Sai, phủ định của ``cân'' là ``không cân'', không phải ``không vuông''. c) Đúng. d) Đúng, phủ định của ``vô tỉ'' là ``hữu tỉ'' (số thực chỉ là hữu tỉ hoặc vô tỉ).}
\end{ex}
